\section{模型优缺点分析}

\subsection{优点}

第一，本文严格识别了两个数据集的重复录音结构，监督学习中按受试者分组交叉验证，并以 subject-level 指标作为主结果，降低了样本单位错误带来的性能高估风险。第二，主声学模型排除了 \feature{sex_male}，避免将人口学变量混入声学诊断结论。第三，候选语音生物标志物识别同时结合统计检验、效应量、稳定选择和模型重要性，并以 XGBoost/SHAP 作补充解释一致性检查，使特征解释不依赖单一方法。第四，对缺少真实标签的第四问和第五问，本文采用无监督潜在表型分析，并明确其探索性边界。第五，全流程一致性检查对分组、泄漏和指标一致性进行了系统控制。

\subsection{局限性}

本文也存在明显局限。Dataset 1 受试者数有限且类别不均衡，导致部分模型特异度偏低；补充 XGBoost 虽提高了部分判别指标，但同样表现出 specificity 偏低的问题，不能据此写成临床确诊模型。Dataset 2 虽类别均衡，但样本量较小。两个数据集的特征空间不完全一致，不能简单合并或直接外推。SHAP 解释反映的是模型预测贡献，不等同于临床因果机制。第四问和第五问缺少真实临床症状标签，因此聚类结果只能称为潜在语音表型，不能直接对应运动型/非运动型或六类临床症状。聚类表型仍需要真实临床标签、纵向随访和外部样本验证。
